\section{Introducción}

Las plataformas musicales generan metadatos y descriptores de audio que permiten caracterizar
catálogos a una escala difícil de abordar manualmente. En este contexto, la inteligencia musical
puede ayudar a describir segmentos, detectar patrones y priorizar preguntas de negocio. Sin
embargo, la disponibilidad de atributos no implica que la popularidad pueda explicarse solo por la
estructura sonora: la exposición, la promoción, la recencia y el reconocimiento previo también
intervienen y no están observados en este dataset.

El proyecto analiza la relación entre atributos de canciones y \code{popularity}. Su alcance actual
incluye comprensión del problema, evaluación de la fuente, preparación, EDA y recomendaciones para
modelamiento. No incluye todavía un modelo entrenado ni una evaluación predictiva fuera de muestra.

Conviene distinguir tres conceptos:

\begin{itemize}
\item \textbf{Asociación:} dos variables cambian de forma relacionada en los datos observados. Pearson y
Spearman cuantifican formas específicas de asociación.
\item \textbf{Predicción:} un procedimiento estima valores para observaciones no usadas al ajustarlo. Requiere
particiones y métricas fuera de muestra.
\item \textbf{Causalidad:} un cambio deliberado en una variable produciría un cambio en otra. Exige un diseño
causal, supuestos o evidencia que este estudio observacional no proporciona.
\end{itemize}

Por tanto, frases como ``mayor sonoridad causa popularidad'' o ``la música explícita tiene éxito'' no
están respaldadas. El lenguaje correcto describe patrones de la muestra y formula hipótesis para
validaciones futuras.
